\chapter{NumPy si Pandas}
\label{chapter:numpy-pandas}

Pentru inceput, un lucru destul de simplu, dar care conteaza mai mult decat pare la
prima vedere: cum tine Python datele in memorie, si de ce, la un moment
dat, un simplu \pycode{list} nu mai e de ajuns.

\section{De ce NumPy}

Cand scrii \pycode{v = [1, 2, 3, 4]} in Python, nu obtii un bloc de numere
asezate frumos unul langa altul in memorie, cum ai crede. Obtii o lista de
referinte, fiecare aratand catre cate un obiect Python separat, cu tot
overhead-ul lui. E flexibil, dar e risipa cand ai milioane de numere de
acelasi fel.

\pycode{numpy.ndarray} rezolva exact asta: tine date brute, de un singur
tip, intr-un bloc continuu de memorie. Dar motivul pentru care merita sa
inveti bine NumPy nu e viteza in sine, e ca te obliga sa gandesti altfel.

In loc sa ii spui calculatorului, pas cu pas, cum sa parcurga fiecare
element, ii spui ce vrei sa obtii, pe tot vectorul deodata. Asta se cheama
\textit{vectorizare}, regula asta o sa te salveze de multe
ori: daca te trezesti ca scrii un \pycode{for} peste un array NumPy,
aproape sigur exista deja o operatie care face exact ce vrei tu, mai
rapid si mai clar.

\begin{minted}{python}
import numpy as np

v = np.array([1, 2, 3, 4])
v * 2      # array([2, 4, 6, 8]), fara bucla
v.mean()   # 2.5
v[v > 2]   # array([3, 4]), indexare booleana
\end{minted}

Nu trebuie sa memorezi functiile astea! Nici la olimpiada nu se
asteapta nimeni sa le stie pe de rost, ai voie cu documentatia acolo.
O sa vezi singur, din exercitii, care le folosesti des si care nu.

\section{Ce face v.mean() de fapt}

Aici e partea care conteaza cu adevarat pentru noi. \pycode{mean()} nu e
magie, e o suma impartita la un numar, atat:

\begin{minted}{python}
def medie(v):
    s = 0
    for x in v:
        s += x
    return s / len(v)
\end{minted}

Asta e tot ce face NumPy pe dinauntru. Doar ca bucla ruleaza intr-un cod
C, peste memorie continua, nu in Python peste obiecte separate, de-aia e
de zeci de ori mai rapit, nu pentru ca ar calcula altceva.
Vezi \hyperref[problem:statistici-vectorizate]{exercitiul} de
mai jos pentru masuratoarea reala.

Ideea se repeta la aproape orice functie din NumPy sau Pandas:
\pycode{sum}, \pycode{std}, \pycode{groupby().mean()}, toate sunt formule
simple, scrise o singura data, corect si rapid, ca sa nu le rescrii tu de
fiecare data. Merita sa stii sa le scrii si de mana, macar o data. Pe asta
ne concentram in cursul asta, nu doar pe ce functie exista si cum se
apeleaza.

\section{Complexitate: aceeasi, dar...}

Amandoua variantele sunt $\mathcal{O}(n)$: fiecare element e vizitat o
singura data, fie in bucla scrisa de mana, fie pe dinauntrul lui
\pycode{v.mean()}. Din punctul de vedere al complexitatii, sunt identice.

\begin{minipage}[t]{0.48\linewidth}
\textbf{Bucla de mana}, $\mathcal{O}(n)$:
\begin{minted}{python}
s = 0
for x in v:
    s += x
medie = s / len(v)
\end{minted}
\end{minipage}
\hfill
\begin{minipage}[t]{0.48\linewidth}
\textbf{NumPy}, $\mathcal{O}(n)$:
\begin{minted}{python}
medie = v.mean()
\end{minted}
\end{minipage}

Atunci de unde diferenta de viteza pe care o masuram mai jos?
Complexitatea iti spune cum creste timpul de executie o data cu $n$, nu
cat dureaza efectiv un pas. Fiecare pas din bucla Python inseamna:
interpretorul citeste o instructiune, verifica tipul lui \pycode{x}, face
un apel de functie pentru \pycode{+=}, etc. Fiecare pas din
\pycode{v.mean()} inseamna o singura adunare in cod masina, pe date deja
asezate in memorie. Constanta ascunsa in spatele lui $\mathcal{O}(n)$ e
mult mai mica la NumPy, chiar daca forma curbei e aceeasi.

De-aia complexitatea nu e de ajuns: iti spune daca un algoritm se
degradeaza bine cand creste $n$, dar nu iti spune care varianta e mai
rapida la un $n$ fix. Pentru asta, masori direct.

\section{Cum masuram timpul de executie}

Am zis mai sus ca varianta vectorizata e "de zeci de ori mai rapida", dar
nu ti-am aratat cum masor asta de fapt. Nu e nimic ascuns: retii timpul
inainte, retii timpul dupa, faci diferenta.

\begin{minted}{python}
import time

start = time.perf_counter()
medie(v)
timp = time.perf_counter() - start

print(timp)   # cat a durat, in secunde
\end{minted}

\pycode{time.perf_counter()} iti da un ceas cat mai precis, facut special
pentru masuratori de genul asta (spre deosebire de \pycode{time.time()},
care e legat de ceasul sistemului de operare si poate sari inainte sau
inapoi, de exemplu la sincronizare de retea). Diferenta dintre cele doua
momente e timpul de rulare, atat.

Ca sa compari doua variante, faci exact acelasi lucru de doua ori si
raportezi una la alta:

\begin{minted}{python}
start = time.perf_counter()
medie(v)
timp_bucla = time.perf_counter() - start

start = time.perf_counter()
v.mean()
timp_numpy = time.perf_counter() - start

print(f"bucla:  {timp_bucla:.4f}s")
print(f"numpy:  {timp_numpy:.4f}s")
print(f"de {timp_bucla / timp_numpy:.0f}x mai rapid")
\end{minted}

Daca lucrezi intr-un notebook Jupyter, ca cele din exercitii, exista si o
comanda mai scurta, \texttt{\%timeit}, care repeta automat masuratoarea de
mai multe ori si iti da o medie mai stabila decat o singura masuratoare
de mana. E convenabila, dar functioneaza doar in Jupyter, nu si intr-un
script \pycode{.py} obisnuit, de-aia merita sa stii si varianta cu
\pycode{time.perf_counter()}: e cea pe care o poti pune oriunde.

\section{De ce Pandas}

Un \pycode{DataFrame} e un tabel: coloane tipate, fiecare coloana fiind pe
dedesubt tot un \pycode{ndarray}, si randuri indexate. E unealta pentru
date tabelare, citire CSV, filtrare, grupare, agregare, pe care o sa o
folosim la orice set de date cu care lucram, inainte sa ajungem macar la
vreun model.

\begin{minted}{python}
import pandas as pd

df = pd.read_csv("date.csv")
df[df["scor"] > 50]                     # filtrare
df.groupby("categorie")["scor"].mean()  # agregare
\end{minted}
